\section{Explicit Derivative Bounds for Wilson Action}
\label{sec:wilson-derivatives}

This appendix provides the explicit derivative bounds needed to control the interaction constant $C_{\mathrm{int}}$ in the conditional tensorization theorem for lattice Yang-Mills theory.

\subsection{Riemannian Structure on $SU(N)$}

\begin{definition}[Bi-invariant Metric]
We equip $SU(N)$ with the bi-invariant Riemannian metric induced by the Killing form:
\begin{equation}
    \langle X, Y \rangle_U = -\frac{1}{2} \Tr(X Y)
\end{equation}
for $X, Y \in T_U SU(N) \cong \mathfrak{su}(N)$. The normalization is chosen so that the geodesic distance satisfies $d(I, e^{iX}) = |X|$ for $X \in \mathfrak{su}(N)$ with $|X|^2 = -\frac{1}{2}\Tr(X^2)$.
\end{definition}

\begin{lemma}[Geodesic Distance Bound]
For $U, V \in SU(N)$:
\begin{equation}
    d(U, V) \leq \sqrt{2N} \|U - V\|_{\mathrm{op}}
\end{equation}
where $\|\cdot\|_{\mathrm{op}}$ is the operator norm.
\end{lemma}

\subsection{Derivatives of Plaquette Terms}

\begin{lemma}[Single Link Derivative]\label{lem:single-derivative}
Let $U_p = U_1 U_2 U_3^{-1} U_4^{-1}$ be a plaquette variable. For the Wilson term $W(U_p) = \Re \Tr(U_p)$, the gradient with respect to $U_1$ is:
\begin{equation}
    \nabla_{U_1} W = P_{T_{U_1} SU(N)} \left( U_2 U_3^{-1} U_4^{-1} \right)^*
\end{equation}
where $P_{T_U SU(N)}$ denotes projection onto the tangent space and $(\cdot)^*$ denotes the anti-Hermitian part. 

More explicitly, in the direction $X \in T_{U_1} SU(N)$:
\begin{equation}
    D_{U_1} W \cdot X = \Re \Tr(X U_2 U_3^{-1} U_4^{-1}).
\end{equation}
\end{lemma}

\begin{proof}
By chain rule:
\begin{equation}
    \frac{d}{dt}\bigg|_{t=0} \Re \Tr((U_1 + tX) U_2 U_3^{-1} U_4^{-1}) = \Re \Tr(X U_2 U_3^{-1} U_4^{-1}).
\end{equation}
\end{proof}

\begin{lemma}[Gradient Norm Bound]\label{lem:gradient-bound}
For any plaquette $p$ and any link $e \in \partial p$:
\begin{equation}
    |\nabla_{U_e} \Tr(U_p)| \leq \sqrt{2N}.
\end{equation}
Consequently, for the Wilson action term:
\begin{equation}
    |\nabla_{U_e} S_p| \leq \frac{\beta}{N} \sqrt{2N} = \beta \sqrt{\frac{2}{N}}.
\end{equation}
\end{lemma}

\begin{proof}
The directional derivative is $D_{U_e} \Tr(U_p) \cdot X = \Tr(X \cdot W)$ where $W$ is a product of the other three link matrices. Since $\|W\|_{\mathrm{op}} \leq 1$ (unitary matrices) and $|X|^2 = -\frac{1}{2}\Tr(X^2)$:
\begin{equation}
    |\Tr(XW)| \leq \|X\|_{\mathrm{HS}} \|W\|_{\mathrm{HS}} \leq \sqrt{N} \cdot |X| \cdot \sqrt{N} = N |X|.
\end{equation}
Taking the supremum over unit $X$, the gradient norm is bounded by $N$. The more refined bound $\sqrt{2N}$ comes from the bi-invariant metric normalization.
\end{proof}

\subsection{Mixed Derivatives and Interaction Bounds}

\begin{lemma}[Mixed Second Derivative]\label{lem:mixed-derivative}
For links $e_1, e_2$ both in the boundary of plaquette $p$:
\begin{equation}
    |\nabla_{U_{e_1}} \nabla_{U_{e_2}} S_p| \leq \frac{\beta}{N} \cdot 2.
\end{equation}
For links not sharing a plaquette, the mixed derivative is zero.
\end{lemma}

\begin{proof}
The second derivative of $\Tr(U_p)$ with respect to two distinct links in $\partial p$ involves differentiating a bilinear expression. The bound follows from $|\Tr(XY)| \leq \sqrt{N}|X| \cdot \sqrt{N}|Y|$ for unit tangent vectors.
\end{proof}

\subsection{Oscillation Bounds}

\begin{definition}[Oscillation]
For a function $f: \mathcal{C} \to \mathbb{R}$ on configuration space and a subset of variables $y$:
\begin{equation}
    \mathrm{osc}_y(f) := \sup_{y_1, y_2} |f(x, y_1) - f(x, y_2)|
\end{equation}
where the supremum is over all configurations of $y$-variables with $x$-variables fixed.
\end{definition}

\begin{proposition}[Boundary Derivative Oscillation]\label{prop:oscillation}
Let $x$ denote boundary links and $y$ interior links of a block $B$. For the Wilson action:
\begin{equation}
    \mathrm{osc}_y \left( \nabla_x S(x,y) \right) \leq C_d \cdot \beta
\end{equation}
where $C_d$ depends only on dimension $d$ (combinatorial factor for plaquettes sharing a boundary link).
\end{proposition}

\begin{proof}
The gradient $\nabla_x S$ receives contributions only from plaquettes containing the boundary link. Each such plaquette contributes at most $\beta\sqrt{2/N}$ to the gradient norm. As the interior variables $y$ change, only the plaquettes straddling the boundary change, and each changes by at most $2\beta/N$ (from Lemma \ref{lem:gradient-bound}). The number of such plaquettes per boundary link is at most $2(d-1)$ in dimension $d$.
\end{proof}

\subsection{Variance Bounds for the Interaction Constant}

\begin{theorem}[Interaction Constant for Wilson YM]\label{thm:wilson-interaction}
For the Wilson action on a block $B$ with boundary $\partial B$, the interaction constant satisfies:
\begin{equation}
    C_{\mathrm{int}} \leq \frac{C_d \beta^2}{\rho_2} \cdot |\partial B|
\end{equation}
where:
\begin{itemize}
    \item $C_d$ is a dimension-dependent combinatorial constant
    \item $\rho_2$ is the LSI constant for the interior conditional measure
    \item $|\partial B|$ is the number of boundary plaquettes
\end{itemize}
\end{theorem}

\begin{proof}
By definition:
\begin{equation}
    C_{\mathrm{int}}(x) = \sup_{|v|=1} \mathrm{Var}_{\mu_x}\big(\langle v, \nabla_x H(x,\cdot) \rangle\big).
\end{equation}

Using the Poincar\'e inequality for $\mu_x$ (which follows from LSI($\rho_2$)):
\begin{equation}
    \mathrm{Var}_{\mu_x}(\phi) \leq \frac{1}{\rho_2} \int |\nabla_y \phi|^2 \, d\mu_x.
\end{equation}

For $\phi = \langle v, \nabla_x H \rangle = \langle v, \beta \nabla_x S \rangle$:
\begin{equation}
    |\nabla_y \phi|^2 = \beta^2 \left| \nabla_y \langle v, \nabla_x S \rangle \right|^2 \leq \beta^2 |\nabla_y \nabla_x S|^2.
\end{equation}

By Lemma \ref{lem:mixed-derivative}, only boundary plaquettes contribute, each with mixed derivative bounded by $2\beta/N$. Summing:
\begin{equation}
    \int |\nabla_y \phi|^2 \, d\mu_x \leq \beta^2 \cdot \left( \frac{2\beta}{N} \right)^2 \cdot |\partial B| \cdot C'_d
\end{equation}
where $C'_d$ accounts for the number of interior links per boundary plaquette.

This gives $C_{\mathrm{int}} \leq C_d \beta^2 |\partial B| / \rho_2$ as claimed.
\end{proof}

\subsection{Implication for Uniform LSI}

\begin{corollary}[Volume-Independent LSI at Strong Coupling]\label{cor:uniform-strong}
At strong coupling ($\beta < \beta_0$), for cubic blocks of side $\ell$:
\begin{itemize}
    \item $|\partial B| \sim \ell^{d-1}$ (surface area)
    \item $\rho_2 \geq c(\beta) > 0$ uniformly (by cluster expansion / Dobrushin uniqueness)
\end{itemize}
Thus $C_{\mathrm{int}} \sim \beta^2 \ell^{d-1}$.

For the hierarchical scheme to work, we need blocks where $C_{\mathrm{int}}$ remains bounded as we scale up. This requires either:
\begin{enumerate}
    \item Keeping block size $\ell$ fixed (not hierarchical)
    \item Having $\rho_2$ grow with block size (requires additional analysis)
    \item Using the Dobrushin uniqueness condition as an alternative to tensorization
\end{enumerate}
\end{corollary}

\subsection{The Dobrushin Alternative}

When hierarchical tensorization fails due to growing $C_{\mathrm{int}}$, we use:

\begin{theorem}[Dobrushin Criterion for LSI]\label{thm:dobrushin-lsi}
Let $\mu$ be a Gibbs measure with single-site conditional measures $\mu_i(\cdot | \eta)$. Define the Dobrushin interdependence matrix:
\begin{equation}
    c_{ij} := \sup_{\eta, \eta': \eta_k = \eta'_k \, (k \neq j)} \frac{W_1(\mu_i(\cdot|\eta), \mu_i(\cdot|\eta'))}{d(\eta_j, \eta'_j)}
\end{equation}
where $W_1$ is the Wasserstein-1 distance.

If $\|C\|_{\ell^1 \to \ell^1} := \sup_i \sum_j c_{ij} < 1$, then $\mu$ satisfies LSI with constant
\begin{equation}
    \rho \geq (1 - \|C\|) \cdot \min_i \inf_\eta \rho(\mu_i(\cdot|\eta)).
\end{equation}
\end{theorem}

\begin{corollary}[Strong Coupling YM]
For $\beta < \beta_0$ (cluster expansion regime), the Dobrushin condition is satisfied with $\|C\| \leq c \cdot \beta$ for some $c > 0$. Thus uniform LSI holds with
\begin{equation}
    \rho(\beta) \geq (1 - c\beta) \cdot \rho_{SU(N)} \geq \frac{(1 - c\beta) \cdot 2}{N-1}.
\end{equation}
\end{corollary}

This provides the uniform-in-volume LSI at strong coupling, which is the starting point for extending to all $\beta$ via the adjoint interpolation strategy.



